
\documentclass{article}
\usepackage{colortbl}
\usepackage{makecell}
\usepackage{multirow}
\usepackage{supertabular}

\begin{document}

\newcounter{utterance}

\centering \large Interaction Transcript for game `mmlu\_pro', experiment `mmlu\_pro\_law', episode 466 with reflexion{-}llama{-}3.1{-}8b{-}instruct.
\vspace{24pt}

{ \footnotesize  \setcounter{utterance}{1}
\setlength{\tabcolsep}{0pt}
\begin{supertabular}{c@{$\;$}|p{.15\linewidth}@{}p{.15\linewidth}p{.15\linewidth}p{.15\linewidth}p{.15\linewidth}p{.15\linewidth}}
    \# & \multicolumn{2}{c}{Player} && \multicolumn{2}{c}{Game Master} \\
    \hline

    \theutterance \stepcounter{utterance}  
    & & & \multicolumn{4}{p{0.6\linewidth}}{
        \cellcolor[rgb]{0.9,0.9,0.9}{
            \makecell[{{p{\linewidth}}}]{
                \texttt{\tiny{[P1$\langle$GM]}}
                \texttt{You have attempted this same task before and did not succeed. These are the plans you wrote after those attempts.} \\
\texttt{Trial \#1: New plan: } \\
\\ 
\texttt{Upon reviewing my previous attempt, I noticed that I was able to answer all the questions correctly, but I failed to achieve the target state of the game, which is to answer all the questions correctly within the given time frame. } \\
\\ 
\texttt{My strategy was to answer each question individually, without considering the overall time constraint. I should have taken a more efficient approach, such as grouping similar questions together or using a more systematic method to answer the questions.} \\
\\ 
\texttt{To improve my strategy, I will take the following steps:} \\
\\ 
\texttt{1. Group similar questions together: I will group questions that require similar knowledge or skills, such as questions about the same topic or questions that require the same type of reasoning.} \\
\texttt{2. Use a more systematic method: I will use a more systematic method to answer the questions, such as creating a checklist or a flowchart to help me keep track of the questions and their answers.} \\
\texttt{3. Manage my time more effectively: I will make sure to allocate enough time for each question and group, and I will use a timer to keep track of the time.} \\
\texttt{4. Review and adjust: I will review my progress regularly and adjust my strategy as needed to ensure that I am on track to answer all the questions correctly within the given time frame.} \\
\\ 
\texttt{By taking these steps, I hope to improve my strategy and achieve the target state of the game.} \\
\\ 
\texttt{Trial \#2: New plan: } \\
\\ 
\texttt{Upon reviewing my previous attempt, I noticed that I was able to answer some questions correctly, but I failed to achieve the target state of the game, which is to answer all the questions correctly within the given time frame. } \\
\\ 
\texttt{My strategy was to answer each question individually, without considering the overall time constraint. I should have taken a more efficient approach, such as grouping similar questions together or using a more systematic method to answer the questions.} \\
\\ 
\texttt{To improve my strategy, I will take the following steps:} \\
\\ 
\texttt{1. Group similar questions together: I will group questions that require similar knowledge or skills, such as questions about the same topic or questions that require the same type of reasoning.} \\
\texttt{2. Use a more systematic method: I will use a more systematic method to answer the questions, such as creating a checklist or a flowchart to help me keep track of the questions and their answers.} \\
\texttt{3. Manage my time more effectively: I will make sure to allocate enough time for each question and group, and I will use a timer to keep track of the time.} \\
\texttt{4. Review and adjust: I will review my progress regularly and adjust my strategy as needed to ensure that I am on track to answer all the questions correctly within the given time frame.} \\
\\ 
\texttt{However, I also noticed that I failed to follow the game rules, specifically the format for answering the questions. I will make sure to follow the format correctly in my next attempt.} \\
\\ 
\texttt{To achieve this, I will:} \\
\\ 
\texttt{* Make sure to answer each question in the format "The answer is (X)" where X is the correct letter choice.} \\
\texttt{* Review the game rules and format before starting the game to ensure I understand what is expected of me.} \\
\texttt{* Take my time to answer each question carefully and make sure to follow the format correctly.} \\
\\ 
\texttt{By taking these steps, I hope to improve my strategy and achieve the target state of the game.} \\
\\ 
\\ 
\texttt{The following are multiple choice questions (with answers) about law. Answer only with "The answer is (X)" where X is the correct letter choice.} \\
\\ 
\texttt{Question:} \\
\texttt{What is the judge ad hoc?} \\
\texttt{Options:} \\
\texttt{A. Judge ad hoc is the president of the ICJ} \\
\texttt{B. Judge ad hoc is a temporary judge appointed for a specific period of time} \\
\texttt{C. Judge ad hoc is the judge that each party will always nominate in every contentious case} \\
\texttt{D. Judge ad hoc is the member of the bench of the ICJ with a casting vote} \\
\texttt{E. Judge ad hoc is a judge who is nominated by the parties involved in a contentious case, irrespective of their nationality} \\
\texttt{F. Judge ad hoc is a judge who decides on the admissibility of cases before the ICJ} \\
\texttt{G. Judge ad hoc is a judge appointed by the Security Council of the United Nations} \\
\texttt{H. Judge ad hoc is a surrogate judge, in case a judge is disqualified or passes away} \\
\texttt{I. If a party to a contentious case before the ICJ does not have a national sitting as judge, it is entitled to nominate someone as a judge solely for that case, with the title of judge ad hoc} \\
\texttt{J. N/A} \\
\texttt{The answer is (I)} \\
\\ 
\texttt{Question:} \\
\texttt{Functions of the law include all but which of the following?} \\
\texttt{Options:} \\
\texttt{A. defining the limits of government power} \\
\texttt{B. regulating the use of public spaces} \\
\texttt{C. keeping the peace} \\
\texttt{D. maximizing individual freedom} \\
\texttt{E. maintaining order and stability} \\
\texttt{F. preventing environmental degradation} \\
\texttt{G. providing a basis for compromise} \\
\texttt{H. promoting social justice} \\
\texttt{I. promoting the principles of the free enterprise system} \\
\texttt{J. encouraging economic growth} \\
\texttt{The answer is (I)} \\
\\ 
\texttt{Question:} \\
\texttt{The \_\_\_\_\_\_\_\_ School of jurisprudence postulates that the law is based on what is "correct."} \\
\texttt{Options:} \\
\texttt{A. Legal Pragmatism} \\
\texttt{B. Legal Formalism} \\
\texttt{C. Comparative} \\
\texttt{D. Analytical} \\
\texttt{E. Sociological} \\
\texttt{F. Historical} \\
\texttt{G. Critical Legal Studies} \\
\texttt{H. Realist} \\
\texttt{I. Positivist} \\
\texttt{J. Natural Law} \\
\texttt{The answer is (J)} \\
\\ 
\texttt{Question:} \\
\texttt{Which word best summarizes Weber's explanation of the development of formally rational law?} \\
\texttt{Options:} \\
\texttt{A. Socialism.} \\
\texttt{B. Legitimacy.} \\
\texttt{C. Authority.} \\
\texttt{D. Democracy.} \\
\texttt{E. Bureaucracy.} \\
\texttt{F. Conflict.} \\
\texttt{G. Capitalism.} \\
\texttt{H. Charisma.} \\
\texttt{I. Co{-}operation.} \\
\texttt{J. Tradition.} \\
\texttt{The answer is (G)} \\
\\ 
\texttt{Question:} \\
\texttt{A state has recently enacted a statute prohibiting the disposal of any nuclear wastes within the state. This law does not contravene or conflict with any federal statutes. A man operates a company in the state that is engaged in the disposal of nuclear wastes. Subsequent to the passage of the state statute, the man, not yet aware of the new law, entered into contracts with many out{-}of{-}state firms to dispose of their nuclear wastes in the state. On account of this new law, however, the man will be unable to perform these contracts. Assume that the man has standing to challenge this state law. Which of the following presents his strongest constitutional grounds to challenge the state law prohibiting the disposal of nuclear wastes within the state?} \\
\texttt{Options:} \\
\texttt{A. The second amendment {-} the right to bear arms.} \\
\texttt{B. The due process clause of the Fourteenth Amendment.} \\
\texttt{C. The tenth amendment {-} powers not delegated to the United States by the Constitution.} \\
\texttt{D. The first amendment {-} freedom of speech.} \\
\texttt{E. The privileges and immunities clause of Article IV, Section 2.} \\
\texttt{F. The commerce clause.} \\
\texttt{G. The sixth amendment {-} right to a fair trial.} \\
\texttt{H. The eighth amendment {-} prohibition of cruel and unusual punishment.} \\
\texttt{I. The equal protection clause of the Fourteenth Amendment.} \\
\texttt{J. N/A} \\
\texttt{The answer is (F)} \\
\\ 
\texttt{Question:} \\
\texttt{A bank teller was on trial for embezzlement. The prosecution offered into evidence the bank's daily deposit and tally records for a six month period, to show that total deposits did not reconcile with the total funds received and recorded. The prosecution tried to enter the records as exhibits without certification or the testimony of an official bank custodian of records. Defense counsel objected on the basis of hearsay. Is the court likely to order the evidence inadmissible without further proof of authenticity?} \\
\texttt{Options:} \\
\texttt{A. Yes, because a qualified bank employee must first authenticate them in person or provide a certification in compliance with the rules.} \\
\texttt{B. No, because the records are considered public documents and do not need further authentication.} \\
\texttt{C. No, because the defense did not provide any counter evidence to dispute the records.} \\
\texttt{D. Yes, because the prosecution must also prove the bank teller had access to manipulate these records.} \\
\texttt{E. Yes, because bank records are not included in the hearsay exception for business records.} \\
\texttt{F. No, because the records are self{-}authenticating and do not require additional proof of authenticity.} \\
\texttt{G. No, because business records kept on a regular basis are allowed as an exception to the hearsay rule and no further qualification is required.} \\
\texttt{H. Yes, because the prosecution failed to provide a chain of custody for the records.} \\
\texttt{I. No, because the trial court has wide discretion in allowing documentary evidence that is regularly kept by a business.} \\
\texttt{J. Yes, because the records must be certified by a public official, not just a bank employee.} \\
            }
        }
    }
     \\ \\

    \theutterance \stepcounter{utterance}  
    & \multicolumn{4}{p{0.6\linewidth}}{
        \cellcolor[rgb]{0.561,0.737,0.561}{
            \makecell[{{p{\linewidth}}}]{
                \texttt{\tiny{[P1$\rangle$GM]}}
                \texttt{The answer is (E).} \\
            }
        }
    }
    & & \\ \\

    \theutterance \stepcounter{utterance}  
    & & & \multicolumn{2}{p{0.3\linewidth}}{
        \cellcolor[rgb]{0.9,0.9,0.9}{
            \makecell[{{p{\linewidth}}}]{
                \texttt{\tiny{[GM$|$GM]}}
                \texttt{E} \\
            }
        }
    }
    & & \\ \\

    \theutterance \stepcounter{utterance}  
    & & & \multicolumn{2}{p{0.3\linewidth}}{
        \cellcolor[rgb]{0.9,0.9,0.9}{
            \makecell[{{p{\linewidth}}}]{
                \texttt{\tiny{[GM$|$GM]}}
                \texttt{A} \\
            }
        }
    }
    & & \\ \\

    \theutterance \stepcounter{utterance}  
    & & & \multicolumn{2}{p{0.3\linewidth}}{
        \cellcolor[rgb]{0.9,0.9,0.9}{
            \makecell[{{p{\linewidth}}}]{
                \texttt{\tiny{[GM$|$GM]}}
                \texttt{game\_result = LOSE} \\
            }
        }
    }
    & & \\ \\

    \theutterance \stepcounter{utterance}  
    & & \multicolumn{4}{p{0.6\linewidth}}{
        \cellcolor[rgb]{0.95,0.95,0.95}{
            \makecell[{{p{\linewidth}}}]{
                \texttt{\tiny{[P1$\leftrightarrow$P]}}
                \texttt{Post{-}episode Reflexion (not sent to GM):} \\
\texttt{New plan: } \\
\\ 
\texttt{Upon reviewing my previous attempt, I noticed that I was able to answer some questions correctly, but I failed to achieve the target state of the game, which is to answer all the questions correctly within the given time frame. } \\
\\ 
\texttt{My strategy was to answer each question individually, without considering the overall time constraint. I should have taken a more efficient approach, such as grouping similar questions together or using a more systematic method to answer the questions.} \\
\\ 
\texttt{However, I also noticed that I failed to follow the game rules, specifically the format for answering the questions. I will make sure to follow the format correctly in my next attempt.} \\
\\ 
\texttt{To achieve this, I will:} \\
\\ 
\texttt{* Make sure to answer each question in the format "The answer is (X)" where X is the correct letter choice.} \\
\texttt{* Review the game rules and format before starting the game to ensure I understand what is expected of me.} \\
\texttt{* Take my time to answer each question carefully and make sure to follow the format correctly.} \\
\\ 
\texttt{Additionally, I will also focus on improving my knowledge of law and increasing my speed and accuracy in answering the questions. I will review the questions and answers from my previous attempts and make sure to understand the reasoning behind each correct answer.} \\
\\ 
\texttt{To improve my knowledge of law, I will:} \\
\\ 
\texttt{* Review the game rules and format before starting the game to ensure I understand what is expected of me.} \\
\texttt{* Focus on understanding the key concepts and principles of law, such as the different types of law, the role of the judiciary, and the principles of justice.} \\
\texttt{* Practice answering questions on law to improve my speed and accuracy.} \\
\\ 
\texttt{By taking these steps, I hope to improve my strategy and achieve the target state of the game.} \\
            }
        }
    }
    & \\ \\

\end{supertabular}
}

\end{document}
